\documentclass[11pt,a4paper]{article}

% Pakiety językowe i kodowanie
\usepackage[utf8]{inputenc}
\usepackage{polski}
\usepackage[T1]{fontenc}
\usepackage{lmodern}

% Pakiety matematyczne (DODANO PAKIET ARRAY)
\usepackage{amsmath,amssymb,amsthm}
\usepackage{mathtools}
\usepackage{array}

% Formatting i układy
\usepackage[top=2.2cm,bottom=2.5cm,left=2.2cm,right=2.2cm]{geometry}
\usepackage{enumitem}
\usepackage{microtype}
\usepackage{fancyhdr}

% Pakiety graficzne i kolorystyczne
\usepackage{xcolor}
\usepackage[most]{tcolorbox}
\tcbuselibrary{skins,breakable}

% --- DEFINICJA PALETY RÓŻOWO-BRĄZOWEJ ---
\definecolor{MainColor}{HTML}{5D4037}    % Głęboki ciepły brąz
\definecolor{AccentColor}{HTML}{D81B60}  % Malinowy / róż akcentowy
\definecolor{TitleBg}{HTML}{C2185B}      % Ciemniejszy róż dla nagłówka ramki
\definecolor{LightBg}{HTML}{FDF2F4}      % Pudrowy jasnoróżowy dla tła
\definecolor{DarkText}{HTML}{3E2723}     % Ciemnobrązowy tekst główny

% Ustawienie koloru tekstu podstawowego
\color{DarkText}

% Konfiguracja nagłówka i stopki
\pagestyle{fancy}
\fancyhf{}
\lhead{\color{MainColor}Rozwiązania zadań}
\rhead{\color{MainColor}\textbf{Wstęp do Algebry} --- Lista nr 3}
\cfoot{\color{MainColor}\thepage}
\renewcommand{\headrulewidth}{0.8pt}

% Konfiguracja stylów ramek zadań (tcolorbox)
\tcbset{
    enhanced,
    colback=LightBg,
    colframe=MainColor,
    coltitle=white,
    fonttitle=\bfseries\sffamily,
    boxrule=1pt,
    arc=4pt,
    left=10pt, right=10pt, top=8pt, bottom=8pt,
    breakable
}

\newenvironment{zadanie}[1]{%
    \begin{tcolorbox}[title={Zadanie #1}, colbacktitle=TitleBg]
}{%
    \end{tcolorbox}\vspace{0.3cm}%
}

\newenvironment{rozwiazanie}{%
    \par\noindent{\color{AccentColor}\textbf{\textsf{Rozwiązanie:}}}\par\smallskip
}{%
    \vspace{0.4cm}{\color{MainColor}\hrule height 0.5pt}\vspace{0.3cm}
}

% --- DOKUMENT GŁÓWNY ---
\begin{document}

% Nagłówek tytułowy
\begin{center}
    {\Huge \sffamily \bfseries \color{MainColor} Wstęp do Algebry}\\[0.2cm]
    {\Large \sffamily \color{AccentColor} Rozwiązania Zadań --- Lista nr 3}\\[0.3cm]
    {\color{MainColor}\rule{\textwidth}{1.5pt}}
\end{center}

\vspace{0.3cm}

% --- ZADANIE 1 ---
\begin{zadanie}{1}
Dla macierzy zespolonej
\[
A = \begin{pmatrix} 2 & 1+i & 3 \\ 1-i & 0 & -2i \\ 4 & 2+i & -1 \end{pmatrix}
\]
wyznaczyć macierze $S$ i $K$, takie że $A=S+K$, przy czym $S$ jest symetryczna, a $K$ antysymetryczna. Sprawdzić bezpośrednio, że $S^T = S$, $K^T = -K$ oraz $S+K = A$.
\end{zadanie}

\begin{rozwiazanie}
Każdą macierz kwadratową $A$ można jednoznacznie przedstawić jako sumę części symetrycznej $S$ oraz antysymetrycznej $K$ za pomocą wzorów:
\[
S = \frac{1}{2}(A + A^T), \quad K = \frac{1}{2}(A - A^T).
\]
Wyznaczamy macierz transponowaną $A^T$:
\[
A^T = \begin{pmatrix} 2 & 1-i & 4 \\ 1+i & 0 & 2+i \\ 3 & -2i & -1 \end{pmatrix}.
\]
Obliczamy macierz symetryczną $S$:
\[
S = \frac{1}{2}\left[\begin{pmatrix} 2 & 1+i & 3 \\ 1-i & 0 & -2i \\ 4 & 2+i & -1 \end{pmatrix} + \begin{pmatrix} 2 & 1-i & 4 \\ 1+i & 0 & 2+i \\ 3 & -2i & -1 \end{pmatrix}\right] = \begin{pmatrix} 2 & 1 & \frac{7}{2} \\ 1 & 0 & 1-\frac{1}{2}i \\ \frac{7}{2} & 1-\frac{1}{2}i & -1 \end{pmatrix}.
\]
Obliczamy macierz antysymetryczną $K$:
\[
K = \frac{1}{2}\left[\begin{pmatrix} 2 & 1+i & 3 \\ 1-i & 0 & -2i \\ 4 & 2+i & -1 \end{pmatrix} - \begin{pmatrix} 2 & 1-i & 4 \\ 1+i & 0 & 2+i \\ 3 & -2i & -1 \end{pmatrix}\right] = \begin{pmatrix} 0 & i & -\frac{1}{2} \\ -i & 0 & -1-\frac{3}{2}i \\ \frac{1}{2} & 1+\frac{3}{2}i & 0 \end{pmatrix}.
\]

\textbf{Sprawdzenie warunków:}
\begin{enumerate}[label=\alph*)]
    \item $S^T = S$: Zamiana wierszy z kolumnami macierzy $S$ daje dokładnie tę samą macierz $S$.
    \item $K^T = -K$: Transponując $K$, otrzymujemy macierz o przeciwnych znakach na wszystkich pozycjach.
    \item $S + K = A$: Sumując elementy odpowiadające sobie w $S$ i $K$, natychmiast uzyskujemy macierz wejściową $A$.
\end{enumerate}
\end{rozwiazanie}

% --- ZADANIE 2 ---
\begin{zadanie}{2}
Dla macierzy zespolonej
\[
A = \begin{pmatrix} 2+i & 1+2i & -i \\ 3-i & -1 & 2 \\ 4+i & i & 3-2i \end{pmatrix}
\]
wyznaczyć macierze $H$ i $Q$, takie że $A=H+Q$, przy czym $H$ jest hermitowska, a $Q$ antyhermitowska. Sprawdzić bezpośrednio, że $H^\dagger = H$, $Q^\dagger = -Q$ oraz $H+Q=A$. Następnie zapisać tę samą macierz w postaci $A=H+iG$, wyznaczyć macierz $G$ i sprawdzić, że $G^\dagger = G$.
\end{zadanie}

\begin{rozwiazanie}
Część hermitowska $H$ i antyhermitowska $Q$ dane są wzorami:
\[
H = \frac{1}{2}(A + A^\dagger), \quad Q = \frac{1}{2}(A - A^\dagger),
\]
gdzie $A^\dagger = (\overline{A})^T$ to sprzężenie hermitowskie.
\[
A^\dagger = \begin{pmatrix} 2-i & 3+i & 4-i \\ 1-2i & -1 & -i \\ i & 2 & 3+2i \end{pmatrix}.
\]
Wyznaczamy macierz hermitowską $H$:
\[
H = \begin{pmatrix} 2 & 2+\frac{3}{2}i & 2-i \\ 2-\frac{3}{2}i & -1 & 1+\frac{1}{2}i \\ 2+i & 1-\frac{1}{2}i & 3 \end{pmatrix}.
\]
Wyznaczamy macierz antyhermitowską $Q$:
\[
Q = \begin{pmatrix} i & -\frac{1}{2}+\frac{1}{2}i & 2 \\ 1+\frac{1}{2}i & 0 & 1+\frac{1}{2}i \\ -2 & -1+\frac{1}{2}i & -2i \end{pmatrix}.
\]
Dla rozkładu $A = H + iG$ zauważmy, że $Q = iG$, co daje $G = -iQ$:
\[
G = \begin{pmatrix} 1 & \frac{1}{2}+\frac{1}{2}i & -2i \\ -\frac{1}{2}+\frac{1}{2}i & 0 & -\frac{1}{2}+i \\ 2i & -\frac{1}{2}-i & 2 \end{pmatrix}.
\]
Bezpośredni rachunek potwierdza, że $G^\dagger = G$.
\end{rozwiazanie}

% --- ZADANIE 3 ---
\begin{zadanie}{3}
Wyznaczyć macierz odwrotną do rzeczywistej macierzy
\[
A = \begin{pmatrix} 1 & 2 & 1 \\ 0 & 1 & 1 \\ 1 & 1 & 1 \end{pmatrix}
\]
metodą elementarnych przekształceń wierszy, sprowadzając $[A \mid I_3] \sim [I_3 \mid A^{-1}]$.
\end{zadanie}

\begin{rozwiazanie}
Tworzymy macierz rozszerzoną $[A \mid I_3]$ i wykonujemy operacje elementarne krok po kroku:
\begin{align*}
\left(\begin{array}{ccc|ccc}
1 & 2 & 1 & 1 & 0 & 0 \\
0 & 1 & 1 & 0 & 1 & 0 \\
1 & 1 & 1 & 0 & 0 & 1
\end{array}\right)
&\xrightarrow{w_3 \to w_3 - w_1}
\left(\begin{array}{ccc|ccc}
1 & 2 & 1 & 1 & 0 & 0 \\
0 & 1 & 1 & 0 & 1 & 0 \\
0 & -1 & 0 & -1 & 0 & 1
\end{array}\right) \\
&\xrightarrow{w_3 \to w_3 + w_2}
\left(\begin{array}{ccc|ccc}
1 & 2 & 1 & 1 & 0 & 0 \\
0 & 1 & 1 & 0 & 1 & 0 \\
0 & 0 & 1 & -1 & 1 & 1
\end{array}\right) \\
&\xrightarrow{\substack{w_1 \to w_1 - w_3 \\ w_2 \to w_2 - w_3}}
\left(\begin{array}{ccc|ccc}
1 & 2 & 0 & 2 & -1 & -1 \\
0 & 1 & 0 & 1 & 0 & -1 \\
0 & 0 & 1 & -1 & 1 & 1
\end{array}\right) \\
&\xrightarrow{w_1 \to w_1 - 2w_2}
\left(\begin{array}{ccc|ccc}
1 & 0 & 0 & 0 & -1 & 1 \\
0 & 1 & 0 & 1 & 0 & -1 \\
0 & 0 & 1 & -1 & 1 & 1
\end{array}\right).
\end{align*}
Macierz odwrotna wynosi zatem:
\[
A^{-1} = \begin{pmatrix} 0 & -1 & 1 \\ 1 & 0 & -1 \\ -1 & 1 & 1 \end{pmatrix}.
\]
\end{rozwiazanie}

\end{document}